\documentclass[conference]{IEEEtran}
\usepackage[utf8]{inputenc}
\usepackage[T1]{fontenc}
\usepackage{cite}
\usepackage{amsmath,amssymb}
\usepackage{graphicx}
\usepackage{booktabs}
\usepackage{url}

\title{Verifier-Guided Evaluation of LLM-Generated Network Configurations: a Paired Study with Shared Initial Candidates}

\author{
\IEEEauthorblockN{Autonomous AI Researcher}
\IEEEauthorblockA{CNSM 2026 Agentic AI Researcher Track}
}

\begin{document}

\maketitle

\begin{abstract}
We preregistered and executed a paired confirmatory study (CAND-2026-001) that measured whether a verifier-triggered guarded pipeline (generate \ensuremath{\rightarrow} validate \ensuremath{\rightarrow} repair, <=1 repair) reduces deterministic-invariant violations relative to single-shot initial generation for intent\ensuremath{\rightarrow}configuration NetOps tasks. Using a deterministic synthetic task generator and a deterministic network verifier, we ran 200 preregistered tasks under shared\_initial\_candidate semantics with the hosted model gpt-5-mini (execution/model\_configuration.json records temperature = null/unset). Baseline configurations passed the verifier in 199/200 tasks (99.5\%); the guarded pipeline passed 200/200 (100\%). Paired absolute risk difference (guarded - baseline) = 0.005; contingency counts n11=199, n10=1, n01=0, n00=0; exact two-sided McNemar p = 1.0; paired-bootstrap 95\% percentile CI = [0.00, 0.015] (10,000 resamples, seed 1729). The single guarded improvement arose from one verifier-triggered repair that corrected a multi-step ordering violation. We report preregistration, full execution accounting (201 model calls: 200 shared initial + 1 repair), paired analysis, representative diagnostic artifacts, and bounded limitations grounded in the archived execution bundle.
\end{abstract}

\section{Introduction}
Generative large language models (LLMs) are increasingly proposed to translate high-level operator intent into device configurations for network operations (NetOps). Operational correctness for such workflows requires generated configurations to satisfy deterministic invariants (reachability, policy compliance, workflow ordering, and group-level safety). Prior prototype systems and benchmarks illustrate feasibility but leave open questions about how much deterministic verification and bounded repair alter acceptance rates when the identical initial sample is reused across comparison conditions.

This paper reports a preregistered paired experiment (CAND-2026-001) designed to isolate the marginal effect of a verifier-triggered automated repair on an identical generated candidate. Under shared\_initial\_candidate semantics the same single initial generation is deterministically reused by both baseline and guarded episodes; any paired difference therefore arises solely from the permitted validator-triggered repair. The primary estimand is the paired difference in per-task binary acceptance by a deterministic verifier (paired\_success\_rate\_difference\_guarded\_minus\_baseline). Because shared\_initial\_candidate semantics were enforced, this estimand measures a repair-only effect and does not capture potential benefits from independent guarded first-pass generations, constrained prompts, or multi-sample selection.

Contributions. (1) A preregistered, fully archived paired experiment executing 200 intent\ensuremath{\rightarrow}configuration tasks with a deterministic verifier and a hosted LLM (gpt-5-mini) under shared\_initial\_candidate semantics. (2) Complete execution accounting and deterministic reconciliation showing 200 shared initial generations and 1 repair call (201 model calls total). (3) Artifact-grounded confirmatory analysis reporting baseline pass 199/200 and guarded pass 200/200 with contingency counts (n11=199,n10=1,n01=0,n00=0), paired risk difference 0.005, exact two-sided McNemar p = 1.0, and bootstrap 95\% CI = [0.00,0.015] (10,000 resamples, seed 1729). (4) A localized failure diagnosis for the single repaired pair using archived validator traces. (5) Detailed reproducibility guidance and bounded limitations tied to archived artifacts and reconciliation records.

\section{Related Work}
We conducted a fresh literature investigation and verified records covering intent\ensuremath{\rightarrow}configuration systems, generate--validate--repair patterns, and agentic-AI safety discussions relevant to NetOps. Prior intent-to-configuration prototypes and task-bench evaluations demonstrate that LLMs can synthesize device configurations from operator intents and that verification is widely advocated as a guardrail [1]. Work such as NetConfEval provides task-oriented benchmark evidence that LLMs can facilitate configuration synthesis, but reported evaluations commonly focus on syntactic or task-level correctness metrics rather than verifier-backed invariant enforcement across deterministic topologies [1].

Deterministic verification and formal checking of configuration invariants have a separate, longer literature in networking. Beckett et al. present a general approach to network configuration verification that emphasizes encoding device semantics and network-wide properties into analyzable specifications; such verifier-focused methods clarify the kinds of invariants a downstream oracle must implement to make generator\ensuremath{\rightarrow}validator pipelines meaningful in practice [2]. Our study situates itself at the intersection of these bodies of work: rather than benchmarking raw generative capability, we measure how an explicit verifier gate and a bounded repair step change acceptance rates when the identical initial candidate is reused.

Cross-domain empirical work on generate\ensuremath{\rightarrow}validate\ensuremath{\rightarrow}repair and test-in-the-loop LLM repair in software engineering provides a methodological precedent. Test-driven or test-in-loop program-repair demonstrations show that pairing generation with a deterministic test oracle (failing test \ensuremath{\rightarrow} patch \ensuremath{\rightarrow} regression test) yields reproducible corrections for many classes of faults, and those studies emphasize artifact archiving to permit independent reproduction [4]. Analogously, tool-augmented agent studies and analyses of LLM agents that call external analyzers or solvers argue that coupling models with deterministic tools materially changes task outcomes and creates audit points amenable to reproducible evaluation [6]. However, these demonstrations are primarily in code or program-synthesis domains; adapting their empirical lessons to NetOps requires careful mapping of network-specific invariants (ordering, group constraints, dependency rules) and verifier semantics.

Survey and reporting-guideline work underline evaluation and reporting challenges for LLM studies. TRIPOD-LLM and broad LLM surveys call for rigorous preregistration, honest reporting of sampling and randomness settings, and preservation of execution artifacts to support reproducibility and avoid selective reporting [3],[5]. These calls motivated our preregistration, deterministic task selection, and archival of provider-call accounting and verifier traces. In particular, TRIPOD-style guidance stresses that reporting model sampling parameters is necessary even when fields are unset; execution bundles must record temperature/null semantics and the reuse policy for initial candidates so readers can correctly interpret paired designs [3].

Comparative synthesis. Across the verified records there is a consistent pattern: (a) promising LLM-based intent\ensuremath{\rightarrow}config prototypes and benchmarks exist, (b) generate\ensuremath{\rightarrow}validate\ensuremath{\rightarrow}repair architectures are advocated and empirically effective in adjacent domains, and (c) end-to-end verifier-integrated NetOps experiments with full artifact grounding and preregistered paired estimands are scarce. Our paired shared\_initial\_candidate design directly addresses this methodological gap by (i) enforcing identical initial candidates across baseline and guarded episodes, (ii) using a deterministic verifier whose per-episode validator\_trace fields provide fine-grained evidence about the nature of failures (parsed\_command\_count, violation\_codes, required\_command\_count), and (iii) archiving provider-call accounting and stage\_counts so that the precise degree of model interaction (shared\_initial\_generation=200, repair=1) is auditable.

Limitations in the prior literature that motivate our approach. Many prior works evaluate generation quality via held-out references or human judgement rather than deterministic invariant enforcement; others integrate emulators or dataplane tests but do not archive the exact paired generation semantics needed to isolate repair-only effects. By contrast, the verifier-centric studies and test-in-the-loop program-repair literature demonstrate that a deterministic oracle can make repair effects measurable and reproducible, but these prior empirical results must be adapted to NetOps because network invariants include ordering and group-level constraints that differ from unit tests used in program repair. Our contribution is therefore methodological and empirical: we apply a preregistered, verifier-grounded paired protocol to quantify the marginal repair effect under shared\_initial\_candidate semantics and provide the full artifact bundle needed to reproduce that narrowly scoped estimand.

\section{Methodology}
Preregistration and design freeze. We preregistered study CAND-2026-001 before execution. The preregistration fixed the primary estimand, the deterministic sampling of 200 tasks drawn by deterministic\_netops\_task\_generator\_v5 (rule-based index selection documented in the manifest), generation\_semantics = shared\_initial\_candidate (one sampled initial generation per task deterministically reused), and allowed at most one automated repair call per task. The preregistration specified the confirmatory test (exact two-sided McNemar via exact binomial tail), the primary failed-call handling (complete\_pair\_primary), a sensitivity analysis treating persistent unparsables as failures, and a power rationale aiming to detect \textasciitilde{}10--15 percentage-point paired increases at N=200 under mid-range baseline assumptions (\textasciitilde{}45--65\%). The preregistration artifact is archived (preregistration\_sha256 in execution manifest).

Task bank, difficulty encoding, and representative statistics. Tasks were generated deterministically (deterministic\_netops\_task\_generator\_v5) and targeted high-difficulty workflow patterns (very\_hard and extreme). Each task manifest encodes initial\_state, expected\_final\_state, required\_sequence (ordered command list), allowed\_touched\_fields, workflow\_policies (e.g., at-least-one-up constraints), difficulty metadata (changed\_interface\_count, dependency\_rule\_count, required\_command\_count), and a normalized reference\_answer. In the archived representative sample required\_command\_count spans 3--9 and changed\_interface\_count up to 3; the single baseline failure occurred on an extreme composite task with required\_command\_count = 9, consistent with concentrated failure risk on higher-complexity tasks (see representative\_task entries).

Model configuration and sampling semantics. The declared requested model was gpt-5-mini (execution/model\_configuration.json). That file records temperature = null (unset). We explicitly note that temperature = null/unset is not evidence of deterministic sampling or temperature = 0; nondeterministic hosted-model sampling remains possible. Provider accounting (deterministic\_reconciliation.json) reports stage\_counts = \{shared\_initial\_generation: 200, repair: 1\} and hosted\_model\_call\_event\_count = 201; cache\_miss\_count = 201, cache\_hit\_count = 0 indicates fresh generation events for the shared-initial calls.

Deterministic verifier and scoring. The deterministic verifier deterministic\_netops\_validator\_v5 served as the oracle. It parses candidate configurations, enforces per-device allowed-field constraints, enforces required\_sequence ordering, evaluates group invariants (e.g., final\_group\_constraints), and returns a binary full\_intent\_constraint\_validation score of 1 iff all encoded invariants pass. Per-episode validator\_trace artifacts include validation\_before and validation\_after states, parsed\_command\_count, required\_command\_count, observed\_touched\_field\_count, violation\_codes, and a boolean valid flag. These verifier traces are archived under execution/scoring/ and execution/responses/ and constitute the authoritative evidence for per-episode scoring.

Paired execution semantics and failure handling. Under shared\_initial\_candidate semantics baseline and guarded episodes for each pair begin from the identical sampled candidate; any paired difference therefore arises solely from the permitted verifier-triggered repair. The missingness\_plan specified a deterministic sanitizer for minor formatting issues and a sensitivity analysis treating persistent unparsables as failures. In the archived run there were no persistent unparsables and no preregistered exclusions; terminal\_error\_type fields are null for all episodes in scoring artifacts.

Statistical analysis and exact McNemar implementation. The primary confirmatory test is the exact two-sided McNemar computed from discordant pair counts n10 (guarded=1, baseline=0) and n01 (guarded=0, baseline=1). We implement the exact two-sided McNemar via the exact binomial tail: with n = n10 + n01 and k = n10, the two-sided p-value is computed as
p = 2 * min\{ BinomCDF(k; n, 0.5), 1 - BinomCDF(k - 1; n, 0.5) \}
with conventional handling for k = 0. For our observed discordants (n = 1, k = 1): BinomCDF(1;1,0.5) = 1 and 1 - BinomCDF(0;1,0.5) = 0.5, so p = 2 * min(1, 0.5) = 1.0. Paired-bootstrap percentile confidence intervals for the absolute risk difference were computed with 10,000 resamples and bootstrap\_seed = 1729; implementation parameters and seed are archived in analysis/analysis\_log.jsonl.

\section{Results}
Execution accounting and provider audit. The archived execution manifest records start and end timestamps (UTC) showing the run elapsed \ensuremath{\approx}29 minutes (start 2026-09-01T16:03:05.491039+00:00, end 2026-09-01T16:32:13.298536+00:00). Planned\_episode\_count = 400 and completed\_episode\_count = 400; failed\_episode\_count = 0. Provider-call accounting (deterministic\_reconciliation.json) reports hosted\_model\_call\_event\_count = 201 with stage\_counts = \{shared\_initial\_generation: 200, repair: 1\}. Cache\_miss\_count = 201 and cache\_hit\_count = 0 indicate fresh model generations for all shared-initial calls and a single repair-stage call; these counts match the preregistered shared\_initial\_candidate semantics and the recorded single repair invocation.

Primary confirmatory outcomes. analysis/condition\_summary.csv shows baseline successes = 199 and failures = 1 (baseline success\_rate = 0.995) and guarded successes = 200 and failures = 0 (guarded success\_rate = 1.000). The paired contingency table (analysis/paired\_contingency\_table.csv) records n11 = 199, n10 = 1, n01 = 0, n00 = 0. From these counts the paired absolute risk difference (guarded - baseline) = 0.005. Using the exact two-sided McNemar formulation described in Methods, with n = n10 + n01 = 1 and k = n10 = 1 we compute BinomCDF(1;1,0.5) = 1 and 1 - BinomCDF(0;1,0.5) = 0.5; hence p = 2 \ensuremath{\times} min(1, 0.5) = 1.0. The paired-bootstrap percentile 95\% confidence interval for the absolute risk difference was computed with 10,000 resamples and bootstrap\_seed = 1729 (analysis/analysis\_log.jsonl) and equals [0.00, 0.015]. These numeric artifacts (analysis/condition\_summary.csv, analysis/paired\_contingency\_table.csv, analysis/analysis\_log.jsonl) are archived and reproduce the reported confirmatory values when the paired\_binary\_analysis\_v1 executor is re-run on the archived input\_results\_sha256.

Single repair event and diagnostic traces. The guarded pipeline invoked exactly one automated repair. Archived validator traces for the repaired pair indicate an extreme composite task (difficulty metadata recorded in the task manifest: required\_command\_count = 9, changed\_interface\_count = 3, difficulty.level = "extreme") whose shared initial candidate violated the task's required\_sequence ordering constraint. The verifier trace for that episode documents validation\_before.violation\_codes showing the ordering breach, the shared\_initial\_candidate (the identical candidate reused by baseline and guarded), and validation\_after with validation\_after.valid = true once the bounded repair was applied. The repair prompt/trace and the validator's normalized\_configuration show the localized reordering the repair enacted; after repair the scorer returned score = 1 (scoring\_reason\_code = "VALID\_CONFIGURATION"). These per-episode fields are present in the execution/scoring/* artifacts and constitute the archived evidence that the single observed paired improvement resulted from a verifier-triggered reordering repair that converted a sequenced-dependency failure into a valid configuration under the verifier's invariants.

Missingness, exclusions, and contamination checks. The run recorded no preregistered exclusions and no persistent unparsables; analysis/missingness\_summary.csv reports complete\_pairs = 200 and zeros for all other missingness categories. analysis/contamination\_summary.csv is empty (no flagged contamination). Deterministic\_reconciliation.json indicates pair\_accounting\_consistent = true and all\_deterministic\_consistency\_checks\_passed = true. No episodes have terminal\_error\_type set, and provider-call audit details align with preregistration (200 shared initial generations, 1 repair). These artifacts support the statement that the confirmatory analysis included the full preregistered sample without post-hoc exclusions.

Exploratory diagnostics by difficulty metadata. Representative task manifests archived in execution/task\_manifest.jsonl encode difficulty metadata (changed\_interface\_count, dependency\_rule\_count, required\_command\_count). In the archived representative subset required\_command\_count values span 3--9 and changed\_interface\_count values up to 3; the lone baseline failure is associated with the extreme composite difficulty pattern (highest required\_command\_count in the sample), consistent with the preregistered exploratory hypothesis that multi-device sequencing and dependency complexity concentrate residual failures. analysis/paired\_difference\_figure.svg (archived) visualizes condition rates and the bootstrap distribution used for the CI; those visual diagnostics and the per-task validator\_trace fields (parsed\_command\_count, observed\_touched\_field\_count, violation\_codes) are available for targeted reanalysis to inspect whether other high-difficulty tasks produced near-miss traces that would benefit from alternative repair strategies.

Synthesis of quantitative uncertainty. The realized baseline success rate (199/200 = 99.5\%) places the confirmatory test in a sparse-discordant regime: with only one discordant pair the exact McNemar p-value gives limited directional evidence (p = 1.0), while the paired-bootstrap CI [0.00, 0.015] quantifies the upper bound of plausible absolute improvement attributable to repair in this repair-only estimand. In practical terms, these results show that, for this synthetic verifier-backed task bank and the single-model run, bounded verifier-triggered repair corrected one expressible sequencing fault; they do not imply large marginal gains across broader settings without further experiments (e.g., independent guarded-generation contrasts, multi-model studies, or live dataplane emulation).

\section{Discussion}
We expand the interpretive analysis and diagnostics here while preserving the paper's existing supported content.

Estimand interpretation and scope (explicit). Because the experiment used shared\_initial\_candidate semantics, baseline and guarded episodes for each task began from the identical sampled candidate; consequently the primary estimand intentionally measures the marginal effect of permitting an automated verifier-triggered repair on that identical candidate (a repair-only effect). This design isolates repair mechanics from any potential benefits that could arise if the guarded condition were allowed to sample independently or draw multiple candidates.

Sparse-discordant regime and inferential consequences. The observed data (n11=199, n10=1, n01=0) place the analysis in a sparse-discordant regime, a setting where conventional asymptotic approximations are unreliable and exact discrete tests or resampling are appropriate. Our use of the exact binomial-tail McNemar calculation and a paired-bootstrap percentile CI (10,000 resamples, seed 1729 as archived in analysis/analysis\_log.jsonl) is consistent with best practice for small discordant counts. Practically, the realized baseline success rate (99.5\%) substantially reduced the experiment's sensitivity to detect modest repair-only effects relative to our preregistered power target (designed for mid-range baseline rates \textasciitilde{}45--65\%). The bootstrap CI [0.00, 0.015] quantifies that the data are compatible with effects up to 1.5 percentage points in absolute terms under our paired repair-only estimand; the preregistered design could not reliably detect larger effects because none were present.

Mechanistic diagnostic of the repaired episode. The single repair call corrected a concrete sequencing violation expressible in the task manifest's required\_sequence field (required\_command\_count = 9 in the repaired manifest). Archived validator\_trace fields show validation\_before.violation\_codes listing the ordering breach, parsed\_command\_count equal to the presented commands, and validation\_after.valid = true after the automated reordering repair. These per-episode traces demonstrate: (1) the verifier's capacity to detect ordering invariants encoded in the manifest; (2) the repair module's ability to perform localized, semantics-preserving reorderings that the verifier accepts; and (3) that such faults concentrated in tasks with higher required\_command\_count and extreme difficulty labels in our deterministic task bank. All diagnostic fields cited above are available in execution/scoring/task-00000X-*.json in the archived bundle.

Construct validity of the deterministic verifier. deterministic\_netops\_validator\_v5 enforces a structured subset of operational invariants encoded in the task payloads (ordering, allowed\_touched\_fields, group-level up/down constraints, and explicit per-field prerequisites). This verifier therefore provides a high-fidelity oracle for manifest-encoded invariants but does not emulate dataplane behavior, vendor runtime idiosyncrasies, or external control-plane side-effects. Construct-valid conclusions should therefore be framed strictly in terms of manifest-specified, deterministically verifiable invariants; the experiment does not claim to validate dataplane reachability or vendor CLI execution semantics that lie outside the verifier's encoding.

Design implications for future NetOps evaluations. The paired shared\_initial\_candidate design cleanly isolates repair capability; however, operational questions remain that require alternative designs: (a) independent guarded-first-pass sampling would quantify whether guarded prompting or constrained generation yields a higher baseline of valid samples before repair; (b) multi-sample selection (n-best) could reveal selection-improvement effects; and (c) multi-model comparisons would quantify model-dependent failure patterns. From an experimental-planning perspective our results highlight that preregistered power calculations must account for plausible high baseline performance---if baseline rates turn out near 100\% the discordant-pair approach will often be underpowered to detect small absolute repairs-only gains.

Reproducibility, artifact usage, and recommended reanalyses. All artifacts necessary to reproduce analysis are archived: execution/task\_manifest.jsonl (deterministic task selection and required\_sequence fields), execution/model\_configuration.json (declared model, noting temperature = null/unset), execution/responses/ (shared-initial and per-condition responses), execution/scoring/ (per-episode validator traces), analysis/paired\_contingency\_table.csv and analysis/condition\_summary.csv, deterministic\_reconciliation.json (provider accounting and stage\_counts), and analysis/analysis\_log.jsonl (bootstrap parameters and seed). Re-running the paired\_binary\_analysis\_v1 executor on the archived input\_results\_sha256 reproduces the confirmatory results. For targeted reanalysis, researchers can (i) recompute bootstrap intervals under alternative resampling schemes, (ii) treat the single repaired task as an outlier to study robustness, or (iii) reclassify manifest invariants to evaluate sensitivity of verifier coverage---all without altering original model-call artifacts.

Threats to validity revisited (artifact-grounded). Internal validity: shared\_initial\_candidate semantics intentionally eliminate differences due to separate sampling, but they also prevent measuring effects of guarded prompt variants; model sampling nondeterminism remains possible because temperature was unset (execution/model\_configuration.json records temperature = null). Construct validity: the verifier enforces only encoded invariants; omissions in the manifest (for example, dataplane-dependent invariants) are not tested. External validity: a single hosted model (gpt-5-mini) and a synthetic deterministic task bank limit generalizability; results do not imply identical outcomes for other models, live devices, or telemetry-driven validation. Conclusion validity: with one discordant pair the statistical evidence for a nonzero repair-only effect is weak (p = 1.0), though the bootstrap CI constrains plausible magnitudes.

Operational implications and cautious hypotheses. Artifact-grounded evidence demonstrates that bounded verifier-triggered repairs can fix localized ordering violations that are detectible and expressible in task manifests. From an operational standpoint this suggests a practical pattern: require explicit, machine-checkable manifest invariants for change tasks and pair LLM synthesis with deterministic validation/repair gates before human or automated deployment. We emphasize these are actionable hypotheses supported by the archived diagnostics in this controlled synthetic setting; they are not operational claims that bounded repair will resolve dataplane-level or adversarially induced failures without further evaluation.

Summary of expanded diagnostics. The archived execution bundle provides the evidentiary chain: deterministic task manifests encoding the invariants, the shared initial generated candidates, per-episode validator traces showing before/after states, provider-call accounting confirming shared\_initial\_generation = 200 and repair = 1, and analysis artifacts (contingency table, bootstrap parameters) documenting the statistical regime. Together these artifacts ground the narrower, repair-only conclusion and illuminate concrete design decisions for future, broader evaluations (independent guarded sampling, multi-model studies, live emulation).

\section{Limitations}
\begin{itemize}
\item Synthetic deterministic task bank and verifier: results reflect correctness under the chosen synthetic intent templates and deterministic verifier and do not guarantee similar performance on live heterogeneous production devices or telemetry.
\item Single-model experiment (gpt-5-mini): findings do not demonstrate multi-model generality.
\item Shared\_initial\_candidate semantics: baseline and guarded conditions began from the same initial generation; the study measures verifier/repair effects only and does not evaluate independent first-pass generation contrasts.
\item Limited adversarial/contamination testing: the run did not include systematic adversarial retrieval or poisoned-context attacks; such threats remain an open evaluation axis.
\item Oracle coverage: the deterministic verifier enforces encoded invariants but may omit runtime behaviors and device-specific semantics observable only through live execution; external validity to live deployments is therefore limited.
\end{itemize}

\section{Conclusion}
We executed a preregistered paired experiment measuring the marginal effect of a verifier-triggered repair under shared\_initial\_candidate semantics for LLM-generated network configurations. In 200 synthetic deterministic tasks with gpt-5-mini and deterministic\_netops\_validator\_v5, baseline acceptance was 199/200 and guarded acceptance was 200/200; a single automated repair corrected a multi-device ordering violation. Paired absolute risk difference = 0.005 with paired-bootstrap 95\% CI [0.00, 0.015] (10,000 resamples, seed 1729) and exact two-sided McNemar p = 1.0. Artifact-grounded evidence therefore supports a constrained conclusion: verifier-guarded bounded repair can correct localized ordering faults in this synthetic verifier-backed setting. Broader operational claims require additional experiments: multi-model studies, independent-first-pass guarded semantics, adversarial contamination testing, and live-device or dataplane emulation to capture runtime semantics not modeled by the deterministic verifier.

\section*{Disclosure Statement}
Disclosure: This run implemented preregistered study CAND-2026-001 and was executed autonomously after the autonomy lock. Declared model: gpt-5-mini (execution/model\_configuration.json records temperature = null/unset; null/unset is not evidence of deterministic sampling or temperature = 0). Orchestration adapter: hosted\_netops\_gvr\_v1. Exact initial master prompt archived at provenance/master\_prompt.txt (SHA-256 1872df1e\allowbreak{}1805d2d9\allowbreak{}6940456c\allowbreak{}a016bd66\allowbreak{}5d1d5196\allowbreak{}add77f5a\allowbreak{}cdf1582b\allowbreak{}b39b15ba). Preregistration fixed generation\_semantics = shared\_initial\_candidate (one sampled initial generation per task was deterministically reused by baseline and guarded episodes) and allowed at most one automated repair call per task. Model-call accounting in the archived execution manifest reflects 200 shared initial generation calls and 1 repair call (201 model calls total). The deterministic verifier used was deterministic\_netops\_validator\_v5. Humans selected the broad topic family and constrained the initial master prompt prior to the autonomy lock as permitted by the track; no human manually edited technical manuscript content after the autonomy lock. All provider traces, verifier logs, task manifests, model-configuration records, and analysis artifacts are archived in the execution bundle and indexed by the execution manifest.

\begin{thebibliography}{99}
\bibitem{ref1} Changjie Wang, Mariano Scazzariello, Alireza Farshin, Simone Ferlin, Dejan Kostić, Marco Chiesa, "NetConfEval: Can LLMs Facilitate Network Configuration?", 2024, 10.1145/3656296
\bibitem{ref2} Ryan Beckett, Aarti Gupta, Ratul Mahajan, David Walker, "A General Approach to Network Configuration Verification", 2017, 10.1145/3098822.3098834
\bibitem{ref3} Jack Gallifant, Majid Afshar, Saleem Ameen, Yindalon Aphinyanaphongs, Shan Chen, Giovanni Cacciamani, Dina Demner‐Fushman, Dmitriy Dligach, Roxana Daneshjou, Chrystinne Oliveira Fernandes, Lasse Hyldig Hansen, Adam Landman, Lisa Soleymani Lehmann, Liam G. McCoy, Timothy A. Miller, Amy C. Moreno, Nikolaj Munch, David Restrepo, Guergana Savova, Renato Umeton, Judy Wawira Gichoya, Professor Gary S. Collins, Karel G.M. Moons, Leo Anthony Celi, Danielle S. Bitterman, "The TRIPOD-LLM reporting guideline for studies using large language models", 2025, 10.1038/s41591-024-03425-5
\bibitem{ref4} Yunhe Li, "Test-in-the-loop LLM Repair: Verifiable Automated Program Repair on QuixBugs with a "Failing Test \ensuremath{\rightarrow} Patch \ensuremath{\rightarrow} Regression Test" Loop", 2024, 10.69987/jacs.2024.40206
\bibitem{ref5} Wayne Xin Zhao, Kun Zhou, Junyi Li, Tianyi Tang, Zican Dong, Yupeng Hou, Beichen Zhang, Yingqian Min, Junjie Zhang, Peiyu Liu, Xiaolei Wang, Yifan Du, Yushuo Chen, Yushuo Chen, Zhipeng Chen, Jinhao Jiang, Ruiyang Ren, Yifan Li, Xinyu Tang, Peiyu Liu, Yiwen Hu, Jian‐Yun Nie, Ji-Rong Wen, "A Survey of Large Language Models", 2026, 10.1007/s11704-026-60308-3
\bibitem{ref6} Haoran Dong, "Tool-Augmented Large Language Model Agents for Analysis: A Focused Study", 2026, 10.2139/ssrn.6647800
\end{thebibliography}

\end{document}
